\section{Grobe Flow-Analyse}
\label{sec:flow-rdy}

\subsection{Annahmen}
\label{sec:annahmen-rdy}

\annahme{Zwei Rechnungen, keine Summe.} Wie in Teil~\ref{sec:flow-azn} sind
die Kapazit\"atsrechnung und die Anlegerrechnung Alternativen. Sie verwenden
denselben Mittelzufluss und d\"urfen nicht addiert werden.

\annahme{Die Entschuldungsfrage entf\"allt.} Das Unternehmen hat eine
Nettokasse von \MioINR{\RdyNettokasse} (Abschnitt~\ref{sec:abschluss-rdy}).
Die Frage \enquote{wie schnell kann es tilgen} ist damit leer, und sie
trotzdem zu rechnen erg\"abe eine Seite Arithmetik ohne Entscheidung. An ihre
Stelle tritt die Frage, die dieses Unternehmen tats\"achlich zu beantworten
hat: Tr\"agt der operative Mittelzufluss das eigene Investitionsprogramm?

\annahme{Freier Mittelzufluss.} Operativer Mittelzufluss abz\"uglich der
Investitionen in Sachanlagen und immaterielle Verm\"ogenswerte, zuz\"uglich
der Erl\"ose aus deren Abgang. Eine Leasingtilgung wird nicht gesondert
abgezogen, weil der Abschluss sie nicht in derselben Gliederung ausweist wie
der von AstraZeneca; die beiden Reihen sind daher in dieser Hinsicht nicht
exakt gleich definiert, und das ist der Grund, warum dieses Dokument die
beiden Titel nicht gegeneinander rechnet.

\annahme{Konstantes Niveau, kein Wachstum.} Wie in Teil~\ref{sec:flow-azn}.
Was der Kurs an Wachstum voraussetzt, wird umgekehrt ausgerechnet.

\annahme{Drei Niveaus.} Die letzten zw\"olf Monate bis zum 30.~Juni 2026
(\INRje{\RdyFcfLtm}), der Median der neun Jahre (\INRje{\RdyFcfMedian}) und
das letzte volle Gesch\"aftsjahr (\INRje{\RdyFcfVoll}).

\annahme{Aktienzahl.} Gerechnet wird durchg\"angig mit
\Stueck{\RdyAktienAusgegeben} Aktien, der Zahl zum 31.~M\"arz 2026. F\"ur die
fr\"uhen Jahre der Reihe ist das eine Vereinfachung; die Aktienzahl hat sich
au{\ss}erhalb des Splits nur wenig ver\"andert, und der Split ist im Abschluss
selbst r\"uckwirkend ber\"ucksichtigt.

\annahme{W\"ahrung.} Die gesamte Rechnung l\"auft in Rupien. Der Kurs des
Hinterlegungsscheins wird einmal umgerechnet, zum Referenzkurs der
Europ\"aischen Zentralbank vom \Datum{\RdyKursTag}
(\Kurswert{\Wechselkurs}~INR je USD).\quelleWeb{Europ\"aische Zentralbank, Euro-Referenzkurse}{quelle:ezb}{17.09.2026}\quelleMerken{s33webxeuropxxaischexzentralbankxxeuroxreferenzkursexquellexezb}
Wer in Dollar denkt, tr\"agt zus\"atzlich das W\"ahrungsrisiko der Rupie; es
ist hier nicht bewertet.

\annahme{Endwert gleich Buchwert je Aktie}, also
\INRje{\RdyBuchwertJeAktieX}. Zwei weitere Endwertannahmen werden daneben
gerechnet, aus demselben Grund wie in Teil~\ref{sec:flow-azn}.

\annahme{Steuern und Horizont.} Wie in Abschnitt~\ref{sec:annahmen-azn}:
nach Konzernsteuern, vor Anlegersteuern, Horizonte von f\"unf, zehn und
f\"unfzehn Jahren, H\"urde \Pct{\Huerde}. Die indische Quellensteuer auf
Dividenden an Auslandsanleger ist nicht ber\"ucksichtigt; sie w\"urde die
Renditen weiter senken.

\subsection{Tr\"agt der Mittelzufluss das Investitionsprogramm?}
\label{sec:kapazitaet-rdy}

\begin{table}[htbp]
\centering
\caption{Finanzierungsrechnung FY2026, Mio.\,INR. Eine Kapazit\"atsrechnung,
keine Prognose.}
\label{tab:rdy-kapazitaet}
\begin{tabular}{lr}
\toprule
Position & Betrag\\
\midrule
\tabellenkoerper{data/flow_rdy_kapazitaet}
\bottomrule
\end{tabular}
\end{table}

\bk{Die Antwort lautet ja, und sie ist knapper, als die Nettokasse vermuten
l\"asst. Nach Investitionen und Dividende bleibt ein \"Uberschuss, der etwa
der H\"ohe der Nettokasse entspricht. Das Unternehmen finanziert sein
Programm also aus einem Jahresertrag, nicht aus Reserven. Die Kehrseite: Im
Quartal nach dem Stichtag war der operative Mittelzufluss negativ
(Abschnitt~\ref{sec:strategie-rdy}). Ein einzelnes Quartal ist keine Reihe;
zwei aufeinanderfolgende w\"aren eine.}

\FloatBarrier

\subsection{Die Rechnung des Anlegers}
\label{sec:anleger-rdy}

\begin{table}[htbp]
\centering
\caption{Interner Zinsfu{\ss} in Prozent pro Jahr beim heutigen Kurs von
\INRje{\RdyKursRupien} je Aktie, Endwert gleich Buchwert je Aktie}
\label{tab:rdy-irr}
\begin{tabular}{lrrr}
\toprule
Niveau des Mittelzuflusses & 5 Jahre & 10 Jahre & 15 Jahre\\
\midrule
\tabellenkoerper{data/flow_rdy_irr}
\bottomrule
\end{tabular}
\end{table}

\begin{table}[htbp]
\centering
\caption{Einstiegsschwelle in INR je Aktie bei einer H\"urde von
\Pct{\Huerde}, Endwert gleich Buchwert je Aktie}
\label{tab:rdy-schwelle}
\begin{tabular}{lrrrr}
\toprule
Niveau & je Aktie & 5 Jahre & 10 Jahre & 15 Jahre\\
\midrule
\tabellenkoerper{data/flow_rdy_schwelle}
\bottomrule
\end{tabular}
\end{table}

\begin{table}[htbp]
\centering
\caption{Dieselbe Schwelle unter der zweiten Endwertannahme: die Aktie notiert
am Ende des Horizonts nominal zum heutigen Kurs}
\label{tab:rdy-schwelle-variante}
\begin{tabular}{lrrrr}
\toprule
Niveau & je Aktie & 5 Jahre & 10 Jahre & 15 Jahre\\
\midrule
\tabellenkoerper{data/flow_rdy_schwellevariante}
\bottomrule
\end{tabular}
\end{table}

\FloatBarrier

\begin{table}[htbp]
\centering
\caption{Umkehrung eins: der Endwert je Aktie, den der heutige Kurs bereits
voraussetzt, in INR und als Vielfaches des Buchwerts}
\label{tab:rdy-endwert}
\begin{tabular}{lrrr}
\toprule
Niveau & 5 Jahre & 10 Jahre & 15 Jahre\\
\midrule
\tabellenkoerper{data/flow_rdy_endwert}
\bottomrule
\end{tabular}
\end{table}

\begin{table}[htbp]
\centering
\caption{Umkehrung zwei: gefordertes j\"ahrliches Wachstum in Prozent bei
konstantem Bewertungsvielfachen}
\label{tab:rdy-vielfach}
\begin{tabular}{lrrr}
\toprule
Niveau & 5 Jahre & 10 Jahre & 15 Jahre\\
\midrule
\tabellenkoerper{data/flow_rdy_vielfach}
\bottomrule
\end{tabular}
\end{table}

\FloatBarrier

\subsection{Was die Tabellen zusammen sagen}

\bk{Wie bei AstraZeneca ist die Rangfolge \"uber alle Horizonte stabil und
wird die H\"urde in keiner Zelle der ersten Tabelle erreicht. Der Unterschied
liegt in der zweiten Umkehrung, und er ist gro{\ss}.}

\bk{Bei konstantem Bewertungsvielfachen verlangt der heutige Kurs
\Pct{\RdyWachstumVielfachZehn} Wachstum des freien Mittelzuflusses pro Jahr.
Ob das viel ist, h\"angt davon ab, welche Vergangenheit man danebenlegt,
und genau hier liegt der Fall dieses Titels: \Pct{\RdyCagr} pro Jahr von
FY2018 an gerechnet, \Pct{\RdyCagrSieben} von FY2019 an. Die eine Rate liegt
\"uber dem Geforderten, die andere ist negativ. Beide stammen aus derselben,
vollst\"andig belegten Reihe.}

\bk{Das ist kein Mangel der Rechnung, sondern eine Eigenschaft des
Unternehmens: Ein Gesch\"aft, dessen Mittelzufluss je Aktie in neun Jahren
zwischen \INRje{\RdyFcfJeAktieXVIII} und dem Vierfachen davon schwankt, hat
keine Trendrate, an der sich ein Preis messen liesse. Teil
\ref{sec:verdikt-rdy} zieht daraus die Folgerung, statt sich eine der beiden
Raten auszusuchen.}
